\documentclass[a4paper, 12pt]{article}

\input{preamble}

\begin{document}
% This sets the default language for the listings package.
\lstset{language=Python}
\title{
\sffamily
\vspace{-1.7cm}
{\huge\textbf{Project 2a -- Alloy Cluster Expansions}} \\[-0.4cm]
}
\author{
    \vspace{-0.4cm}
    \begin{minipage}[t]{0.45\textwidth}
\centering
Linus Brink\\
\href{mailto:brinkl@chalmers.se}{brinkl@chalmers.se}
\end{minipage}%
\hfill
\begin{minipage}[t]{0.45\textwidth}
\centering
Oscar Stommendal\\
\href{mailto:oscarsto@chalmers.se}{oscarsto@chalmers.se}
\end{minipage}
}
\date{\today}
\maketitle
% \begin{center}
%     \textbf{Abstract}
% \end{center}
\vspace{-1.5cm}
\section{Introduction}
This report explores different methods for constructing an alloy cluster expansion model for the Au-Cu alloy, using various levels of physical intuition in the model construction process.
This physical intuition will be instilled through more and more complex priors, ranging from the completely uninformative prior of OLS to a full Bayesian analysis where the prior hyperparameters are sampled as well.
The different cluster expansion models are then applied to the problem of predicting the ground state structure amongst a few different candidates.

\section{Theory}

Density functional theory (DFT) provides the total energy
$E_{\text{structure}}$ of a given atomistic configuration.
For alloy systems it is convenient to work instead with the mixing energy
per atom, which measures the energy relative to the pure elements.
For an Au--Cu structure this is defined as
\begin{equation}
  E_{\text{mix}} =
  \frac{E_{\text{structure}} - n_{\text{Au}} E_{\text{Au}}
        - n_{\text{Cu}} E_{\text{Cu}}}
       {n_{\text{Au}} + n_{\text{Cu}}},
\end{equation}
where $n_{\text{Au}}$ and $n_{\text{Cu}}$ are the numbers of Au and Cu atoms
and $E_{\text{Au}}$ and $E_{\text{Cu}}$ are the energies of pure Au and Cu.
The mixing energy is the relevant quantity when analysing, for example,
phase stability of alloys.

The alloy cluster expansion provides an atomistic model on a perfect lattice.
In this framework the mixing energy is written as a sum over
symmetry–distinct clusters (orbits),
\begin{equation}
  E_{\text{mix}} = J_0 + \sum_{\alpha} m_{\alpha} J_{\alpha},
\end{equation}
where $m_{\alpha}$ is the multiplicity of orbit $\alpha$ and $J_{\alpha}$
the corresponding effective cluster interaction (ECI).
For a set of configurations this expression can be cast in linear form,
\begin{equation}
  \mathbf{E} = \mathbf{X}\mathbf{J},
\end{equation}
where $\mathbf{E}$ collects the mixing energies, $\mathbf{X}$ is the matrix
of cluster vectors and $\mathbf{J}$ is the ECI vector.

\section{Method}



\printbibliography

\end{document}